\begin{agradecimentos}
À professora Maria Helena Cautiero Horta Jardim, pela orientação, pelo incentivo fundamental ao nos convidar para apresentar o projeto original na 13ª Semana de Integração Acadêmica da UFRJ (SIAC) e por aceitar o desafio de coorientar este TCC desde a sua concepção como trabalho de disciplina.

À Luziane Ferreira de Mendonça, por aceitar o papel de orientação, pela análise do projeto inicial e definição dos rumos do Trabalho de Conclusão de Curso.

Aos colegas João Matheus Nascimento Gonçalves e Vitor Lucio Giorgio Cardoso de Carvalho, pela parceria e colaboração intelectual durante a elaboração do projeto original que serviu de alicerce para este trabalho.

À Universidade Federal do Rio de Janeiro (UFRJ), pelo ambiente acadêmico e pelas oportunidades de crescimento proporcionadas ao longo da graduação.

À Secretaria Municipal da Saúde do Rio de Janeiro, pela oportunidade de realizar o estágio que me permitiu aperfeiçoar metodologias ágeis de trabalho.
\end{agradecimentos}